//---------------------------------------
// Tác giả macro: NGUYỄN THANH PHONG
// Chức năng: Kỹ thuật vẽ hình học phẳng cơ bản (TikZ)
// Gói phụ thuộc: \usetikzlibrary{calc, intersections, angles, quotes}
//---------------------------------------
var box = new UniversalInputDialog();
box.setWindowTitle("Kỹ thuật vẽ Hình Học Phẳng");

// Danh sách hiển thị trong ComboBox
var dsHinhHoc = [
"1. Đánh dấu góc",
"2. Đặt tên góc (có text)",
"3. Vòng lặp FOREACH vẽ điểm",
"4. Vẽ tam giác khi biết 3 cạnh",
"5. Vẽ tam giác vuông",
"6. Vẽ tam giác đều",
"7. Vẽ tam giác cân",
"8. Phép quay và vị tự",
"9. Giao điểm của hai đường thẳng",
"10. Giao điểm 2 đường tròn (lệnh let)",
"11. Giao điểm của đường thẳng và đường tròn",
"12. Xác định 2 tiếp điểm (lệnh let)"
];

box.add(dsHinhHoc, "Chọn đoạn code cần chèn:", "LuaChonHH");

// Mảng chứa code TikZ tương ứng
var dataHinhHoc = [
// 1. Đánh dấu góc
[
"\t\\draw pic[draw,angle radius=12mm]{angle=C--B--A}; %Theo chiều dương, tùy chọn double,fill=yellow!50,-stealth,dashed"
].join("\n"),

// 2. Đặt tên góc
[
"\t\\draw pic[\"$\\alpha$\",-stealth,angle radius=8mm]{angle=A--C--B}; %Thêm dấu ngoặc kép"
].join("\n"),

// 3. FOREACH
[
"\t\\foreach \\x/\\g in {A/90,B/180,C/-45,D/0}",
"\t\t\t\\fill[black] \t(\\x) circle (1pt)",
"\t\t\t($(\\g:3mm)+(\\x)$) node {$\\x$};"
].join("\n"),

// 4. Vẽ tam giác khi biết 3 cạnh
[
"\t\\def\\a{3}",
"\t\\def\\b{2}",
"\t\\def\\c{4}",
"\t\\path \t(0:0) coordinate (B)",
"\t\t\t++(0:\\a) coordinate (C);",
"\t\\path[name path=c1]  (B) circle (\\c);",
"\t\\path[name path=c2]  (C) circle (\\b);",
"\t\\path[name intersections={of=c1 and c2,by={A,D}}];",
"\t\\pgfresetboundingbox %Co khung hình ",
"\t\\draw (A)--(B)--(C)--cycle;",
"\t\\foreach \\x/ \\goc in {A/135,B/-135,C/-45} ",
"\t\t\t\\fill (\\x) circle (1pt)",
"\t\t\t($(\\x)+(\\goc:3mm)$) node {$\\x$};"
].join("\n"),

// 5. Vẽ tam giác vuông
[
"\t\\def\\r{3}",
"\t\\path \t(180:\\r) coordinate (B)",
"\t\t\t(0:\\r) coordinate (C)\t",
"\t\t\t(130:\\r) coordinate (A);",
"\t\\draw (A)--(B)--(C)--cycle",
"\t\t\tpic[draw,angle radius=2mm]{right angle=B--A--C};",
"\t\\foreach \\x/ \\goc in {A/135,B/-135,C/-45} ",
"\t\t\t\\fill (\\x) circle (1pt)",
"\t\t\t($(\\x)+(\\goc:3mm)$) node {$\\x$};"
].join("\n"),

// 6. Vẽ tam giác đều
[
"\t\\def\\a{3}",
"\t\\path (0:0) coordinate (B)",
"\t\t\t++(0:\\a) coordinate (C)\t",
"\t\t\t([turn]120:\\a) coordinate (A);",
"\t\\draw (A)--(B)--(C)--cycle;",
"\t\\foreach \\x/ \\goc in {A/135,B/-135,C/-45} ",
"\t\t\t\\fill (\\x) circle (1pt)",
"\t\t\t($(\\x)+(\\goc:3mm)$) node {$\\x$};"
].join("\n"),

// 7. Vẽ tam giác cân
[
"\t\\def\\a{3}",
"\t\\def\\h{3.5}",
"\t\\path \t(0:0) coordinate (B)",
"\t\t\t++(0:\\a) coordinate (C)",
"\t\t\t($(B)!0.5!(C)$) coordinate (H)\t",
"\t\t\t($(H)+(90:\\h)$) coordinate (A);",
"\t\\draw (A)--(B)--(C)--cycle;",
"\t\\foreach \\x/ \\goc in {A/135,B/-135,C/-45} ",
"\t\t\t\\fill (\\x) circle (1pt)",
"\t\t\t($(\\x)+(\\goc:3mm)$) node {$\\x$};"
].join("\n"),

// 8. Phép quay và vị tự
[
"\t($(A)!k!30:(B)$) coordinate (B')"
].join("\n"),

// 9. Giao điểm của hai đường thẳng
[
"\t(intersection of A--B and C--D) coordinate (I)"
].join("\n"),

// 10. Giao điểm thứ 2
[
"\t\\path [name path=c1] (O_1)  let\\p1=($(O_1)-(A)$) in circle ({veclen(\\x1,\\y1)});",
"\t\\path [name path=c2] (O_2) circle (\\r);",
"\t\\path [name intersections={of=c1 and c2,by={C,D}}];"
].join("\n"),

// 11. Giao điểm của đường thẳng và đường tròn
[
"\t\\path \t[name path=c1] (O_1) circle (\\R);",
"\t\\path \t[name path=NAt] (N)--(At);",
"\t\\path \t[name intersections={of=c1 and NAt,by={A}}];"
].join("\n"),

// 12. Xác định 2 tiếp điểm
[
"\t\\path \t[name path=c2] (O_2)  let\\p1=($(O_2)-(M)$) in circle ({veclen(\\x1,\\y1)});",
"\t\\path \t[name path=c3] ($(O_2)!0.5!(E)$) let\\p1=($(O_2)!0.5!(E)-(E)$) in circle ({veclen(\\x1,\\y1)});",
"\t\\path \t[name intersections={of=c2 and c3,by={J,I}}];"
].join("\n")
];

// Khởi chạy hộp thoại
if (box.exec() != null) {
	var result = String(box.get("LuaChonHH"));
	var index = -1;
	
	// Tìm index
	for (var i = 0; i < dsHinhHoc.length; i++) {
		if (dsHinhHoc[i] == result) {
			index = i;
			break;
		}
	}
	
	// Chèn code
	if (index > -1) {
		editor.insertSnippet(dataHinhHoc[index]);
	} else {
		alert("Macro chưa nhận dạng được mục bạn vừa chọn!");
	}
}